\chapter*{Preface: What a Tutor Teaches}
\addcontentsline{toc}{chapter}{Preface: What a Tutor Teaches}

A keyboard layout specifies which movements are possible. A typing tutor specifies which sequences of movements become thinkable, habitual, and eventually invisible. This book is about the second specification, which is easy to overlook because it is rarely stated. It lives in exercises.

\section*{The Method}

Every typing tutor contains a theory of skill, whether or not its authors wrote one down. A lesson that adds one key at a time assumes that skill accumulates by small extensions of mastered territory. A lesson that rewards speed regardless of errors assumes that velocity is the primary quantity. A lesson built from digraphs treats language as a distribution of transitions, while one built from isolated letters treats the keyboard as a set of positions. These commitments can be recovered from the exercises themselves.

The method has three steps.

\begin{enumerate}
\item Treat the exercise as primary evidence and describe it exactly: what is on the page, in what order, repeated how many times.
\item State the generative rule that would reproduce it. Chapter~8 does this for one printed curriculum, whose lessons admit a growing key set $\Kset_i$ and therefore a growing available vocabulary
\[
\Vocab_i=\{\,w\in\Vocab:\Sig(w)\subseteq \Kset_i\,\}.
\]
\item Compare rules across tutors, layouts, and media, so that a difference in output can be traced to a difference in the rule instead of being left as an impression.
\end{enumerate}

\section*{What the Book Does Not Claim}

Four limits are stated at the outset.

The paper tutors of Chapter~8 rest on a small number of photographed pages from two printed sources, and the pages do not settle how the sources were used together. Where the pages are silent, the text marks its statements as inference. Where a line is legible only as show-through from the reverse of a page, it is not quoted as text.

The history in Parts~I to~IV is organized around one question, namely how inscription became a repeatable sensorimotor procedure. It is not a complete history of writing machines, and it omits much that a history of any single device would require. Dates and attributions that could not be checked while drafting carry a \textbf{[Source needed]} flag.

Claims about layouts, Dvorak above all, have been contested for decades. The book reports the dispute without settling it. The curricula of Part~VII compare layouts under identical pedagogy precisely so that the comparison does not depend on the disputed historical claims.

The DOS material of Part~VI is specified as a framework and an evidence schedule. Statements about particular programs are to be filled from emulated sessions, manuals, and screenshots, and are flagged until then.

\section*{The Argument in Outline}

Each part removes an assumption that seems obvious. Part~I shows that writing need not preserve ordinary spelling. Part~II shows that a fixed array of keys can stand in for the moving hand, and that institutions, not only machines, trained the body to use it. Part~III shows that a keypress need not make a mark on paper, and that ``typing'' has named several distinct kinds of work. Part~IV shows that a keystroke need not be a single key, and that the arrangement of keys is a variable. Part~V supplies the vocabulary: implicit theories of skill, a taxonomy of tutors, and a learner state $S=(A,V,H,E,T,R)$ that explains why two typists with the same words per minute need different lessons. Part~VI turns to the DOS era, when tutors became legible enough to read as machines. Part~VII presents the exercises themselves, and Part~VIII specifies a tutor that keeps a complete ledger of practice.

The last assumption to go is that an exercise is neutral because its content is trivial. A line such as \texttt{a sad lad asks dad;} is not merely poor prose. It is the visible shape of a restriction on what may be said, and different restrictions produce different temporary Englishes.

\section*{How to Read It}

The chapters are cumulative through Part~V and can be read more freely afterward. Readers with a particular interest in software may begin with Chapter~8, which carries the method in its most concrete form, and return to the earlier history afterward. Mathematical notation is gathered in Appendix~D, and every formal statement is accompanied by a plain-language account. Transcriptions of the exercise pages used as evidence appear in Appendix~C, and Appendix~F lists every source page with its status, so that any claim about a lesson can be checked against the lesson.

Three evidence labels recur. Text that describes what is printed on a page is observation. Text marked \emph{Inference} draws a conclusion from what a page omits. Text marked \emph{Design} describes a proposal, such as a planned curriculum, and must not be read as a report of what any historical tutor did.

The central contention is simple and is tested throughout: what a person can type is decided by the keyboard, but what a person learns to type is decided by the tutor.
